\documentclass{beamer}
\usepackage{risk201}

\setchapter{9}

% Chapter-specific diagram shortcuts
% Decreasing, convex price curve with tangent line (duration / convexity)
\newcommand{\pricecurve}[1]{%
  \begin{tikzpicture}[scale=#1]
    \draw[->, thick] (0,0) -- (6,0);
    \draw[->, thick] (0,0) -- (0,3.6);
    \node at (-0.4,3.0) {\large P};
    \draw[thick, domain=0:5, samples=60] plot (\x, {0.55+2.7*exp(-0.9*\x)});
    % tangent at x0 = 1.2
    \draw[thick] (0.05,{1.474+0.832*1.15}) -- (2.6,{1.474-0.832*1.4});
    \draw[thick, dashed] (0,1.474) -- (1.2,1.474);
    \draw[thick, dashed] (0,1.0) -- (2.0,1.0);
    \draw[thick] (1.2,0.2) -- (1.2,-0.2);
    \draw[thick] (2.0,0.2) -- (2.0,-0.2);
  \end{tikzpicture}%
}


\title{RISK 201: Asset-Liability Management}

\begin{document}
\maketitle

\begin{frame}{Assets and liabilities}

\begin{itemize}
\item Insurance company \textbf{liabilities} consist of future cash flows promised to policyholders (e.g., annuity or pension payments, payouts on settlements after an injury, etc.).
\medskip
\item These liabilities must be backed by \textbf{assets}.
\medskip
\item The best assets to back a liability stream are ones that produce cash flows with the \green{same timing} and in the \green{same amount} as the liabilities. This is called \purple{exact cash flow matching} or \purple{dedication}.
\end{itemize}

\end{frame}

\begin{frame}{Exact cash flow matching}

\begin{itemize}
\item Creating a dedicated portfolio \green{using zero-coupon bonds} is easy $\rightarrow$ buy bonds to match CFs one by one
\medskip
\item When \purple{using bonds with coupons}:
\begin{enumerate}
\item Consider the final liability cash flow first.
\item Find the amount of the asset needed to match the last liability cash flow.
\item Determine the liability cash flows that are left over after the first matching asset is chosen.
\item Consider the longest \purple{remaining} cash flow; repeat steps 2--4 until all liability cash flows are matched.
\end{enumerate}
\end{itemize}

\end{frame}

\begin{frame}{Example 9.1}

A company has an obligation to pay \$10,000 at the end of each year for 3 years. Spot rates are 4\%, 5\%, and 5.8\% for terms of 1, 2 and 3 years, respectively.
\begin{enumerate}[a)]
\item Construct a dedicated portfolio to match the cash flows using 1-, 2- and 3-year zero coupon bonds.
\end{enumerate}

\begin{center}
\cfrows{6.6}{0/0, 2.2/1, 4.4/2, 6.6/3}{dgreen/1-year zero:, purple/2-year zero:, red/3-year zero:}
\end{center}

\vfill
\end{frame}

\begin{frame}{Example 9.1, continued}

\begin{enumerate}[a)]
\setcounter{enumi}{1}
\item Construct a dedicated portfolio using the following:
\begin{itemize}
\item Bond 1: 3-year bond with 5\% annual coupons
\item Bond 2: 2-year bond with 2\% annual coupons
\item Bond 3: 1-year zero-coupon bond
\end{itemize}
\end{enumerate}

\begin{center}
\cfrows{6.6}{0/0, 2.2/1, 4.4/2, 6.6/3}{dgreen/3-year bond:, purple/2-year bond:, red/1-year zero:}
\end{center}

\vfill
\end{frame}

\begin{frame}{Interest rate sensitivity}

\begin{itemize}
\item Creating dedicated portfolios is not always possible; assets and liabilities might not be perfectly matched.
\medskip
\item Even if we pick assets whose PV matches the PV of the liabilities, the PVs might drift apart when interest rates change and
\begin{itemize}
\item the value of fixed-income assets, like bonds, changes
\item the (market-consistent) value of the liabilities changes
\end{itemize}
\medskip
\item \underline{Goal}: choose assets so that the PV of the assets moves in unison with the PV of the liabilities \\
$\rightarrow$ protect \purple{surplus} from changes in interest rates
\end{itemize}

\end{frame}

\begin{frame}{Surplus}

\begin{itemize}
\item Define the \purple{surplus}, \purple{$S(i)$}, as the PV of assets minus the PV of liabilities, all valued at rate $i$:
\[
S(i) = \sum_{t=0}^{n} \frac{A_t}{(1+i)^t} - \sum_{t=0}^{n} \frac{L_t}{(1+i)^t}
\]
\item Select a portfolio so that $S(i) > 0$ under certain conditions. This is called \green{immunization}.
\medskip
\item To do so, we need to understand how the PV of cash flows change in response to changes in interest rates.
\end{itemize}

\end{frame}

%======================================================================
\section{Duration}
%======================================================================

\begin{frame}{9.1 Duration}

\begin{itemize}
\item Given annual cash flows $C_t$ and continuously compounded yield $\delta$, the price at time 0 is
\[
P(\delta) = \sum_{t=0}^{n} C_t\, e^{-\delta t}
\]
\end{itemize}

\begin{columns}[c]
\begin{column}{0.35\textwidth}
\begin{itemize}
\item Price curve is usually decreasing and convex.
\end{itemize}
\end{column}
\begin{column}{0.6\textwidth}
\centering
\pricecurve{0.85}
\end{column}
\end{columns}

\end{frame}

\begin{frame}{Approximating the price change}

\begin{itemize}
\item The price change resulting from a change in the continuously compounded yield is
\[
P(\delta + \Delta\delta) = P(\delta) + (\Delta\delta) P'(\delta) + \frac{(\Delta\delta)^2 P''(\delta)}{2!} + \cdots
\]
\item A first-order approximation gives
\[
P(\delta + \Delta\delta) \approx P(\delta) + (\Delta\delta) P'(\delta)
\]
\item Graphically, this is a \purple{linear approximation}, based on the current price, $P(\delta)$, and the \purple{tangent} to the price curve at $\delta$.
\end{itemize}

\end{frame}

\begin{frame}{Percentage change in price}

\begin{itemize}
\item For bonds, we are often interested in the \purple{percentage change in the price} in response to a change in the (continuously compounded) yield $\delta$:
\[
\purple{\%\Delta P(\delta)} = \frac{\Delta P(\delta)}{P(\delta)} = \frac{P(\delta + \Delta\delta) - P(\delta)}{P(\delta)} = (\Delta\delta)\, \green{\frac{P'(\delta)}{P(\delta)}}
\]
\end{itemize}

\end{frame}

\begin{frame}{Duration}

\green{(Macaulay) duration} is the \underline{negative} of the \underline{derivative} of the price function with respect to the continuously compounded yield, \underline{divided} by the price:
\[
\green{\bar{d}} = -\frac{\frac{d}{d\delta} P(\delta)}{P(\delta)}
\]

If we know the cash flows' duration, we can easily approximate the \% change in price from a change in $\delta$:
\[
\purple{\%\Delta P(\delta)} \approx (\Delta\delta)\, \frac{P'(\delta)}{P(\delta)} = -\green{\bar{d}}\,(\Delta\delta)
\]

\end{frame}

\begin{frame}{Measuring the change in yield}

\begin{itemize}
\item The change in the yield is measured in \textcolor{red}{absolute} terms. We often express this in basis points: \\
1 basis point (bp) = 0.01\% \quad (100 bps = 1\%)
\medskip
\item When we say a rate is $x$ bps higher/lower than some other rate, we always mean this \textcolor{red}{additively}.
\medskip
\item \underline{Example 9.2}: Suppose the 30-year continuously compounded yield observed in the market on Apr 1 is 10\%. Determine the yield on April 15 if you are told it is 5 bps higher than the rate on April 1.
\end{itemize}

\vfill
\end{frame}

\begin{frame}{Interpreting (Macaulay) duration}

\[
P(\delta) = \sum_{t=0}^{n} C_t\, e^{-\delta t} \;\Rightarrow\; \frac{d}{d\delta} P(\delta) = \sum_{t=0}^{n} -t \cdot C_t\, e^{-\delta t}
\]
\[
\bar{d} = \frac{\sum_{t=0}^{n} t \cdot C_t\, e^{-\delta t}}{\sum_{t=0}^{n} C_t\, e^{-\delta t}}
= \sum_{t=0}^{n} t \cdot \left(\frac{C_t\, e^{-\delta t}}{\sum_{t=0}^{n} C_t\, e^{-\delta t}}\right)
= \sum_{t=0}^{n} t \cdot \frac{PV_t}{PV_{all}}
\]

\medskip
Macaulay duration is the \purple{weighted average of the times that cash flows occur}, where the weights are the relative contributions of each cash flow to the price.

\end{frame}

\begin{frame}{Observations}

\begin{enumerate}
\item The (Macaulay) duration of an $n$-year zero-coupon bond is $n$.
\medskip
\item A bond with a higher coupon rate has a shorter Macaulay duration because the higher coupons mean more weight is placed on earlier payments and less weight on late payments.
\medskip
\item A bond with a higher yield has a shorter Macaulay duration because the relative contribution of the later cash flows to the price is smaller when the yield is higher, so these cash flows have a smaller weight.
\end{enumerate}

\end{frame}

\begin{frame}{Example 9.3}

A 4-year bond has 6\% annual coupons and is priced at par.
\begin{enumerate}[a)]
\item Find the duration of this bond.
\end{enumerate}

\begin{center}
\renewcommand{\arraystretch}{1.3}
\begin{tabular}{cccc}
\toprule
$t$ & $C_t$ & $C_t e^{-\delta t}$ & $t \cdot C_t e^{-\delta t}$ \\
\midrule
1 & $0.06F$ & $0.06F/(1.06)$ & $1 \cdot 0.06F/(1.06)$ \\
2 & $0.06F$ & $0.06F/(1.06)^2$ & $2 \cdot 0.06F/(1.06)^2$ \\
3 & $0.06F$ & $0.06F/(1.06)^3$ & $3 \cdot 0.06F/(1.06)^3$ \\
4 & $1.06F$ & $1.06F/(1.06)^4$ & $4 \cdot 1.06F/(1.06)^4$ \\
\bottomrule
\end{tabular}
\end{center}

\vfill
\end{frame}

\begin{frame}{Example 9.3, continued}

\begin{enumerate}[a)]
\setcounter{enumi}{1}
\item Approximate the percentage change in price resulting from a 10 bps increase in the continuously compounded yield.
\end{enumerate}

\vfill
\end{frame}

%======================================================================
\section{Modified duration (volatility)}
%======================================================================

\begin{frame}{9.2 Modified duration (volatility)}

\begin{itemize}
\item Yield is usually expressed as an effective rate of interest per period, $i$, not a cont.\ comp.\ rate $\delta$.
\medskip
\item We are normally interested in knowing/predicting the percentage change in price on account of a change in the yield $i$ (not $\delta$).
\medskip
\item \orange{Modified duration} (aka \orange{volatility}) is defined as the negative derivative of the price function with respect to the \orange{effective yield}, divided by the price:
\[
\orange{\bar{v}} = -\frac{\frac{d}{d\orange{i}} P(\orange{i})}{P(\orange{i})}
\]
\end{itemize}

\end{frame}

\begin{frame}{Connection between $\bar{d}$ and $\bar{v}$}

\begin{itemize}
\item With annual effective rate $i$:
\[
P(i) = \sum_{t=0}^{n} C_t\,(1+i)^{-t} \;\Rightarrow\; \frac{d}{di} P(i) = \sum_{t=0}^{n} -t \cdot C_t\,(1+i)^{-t-1}
\]
\[
\orange{\bar{v}} = \frac{\sum_{t=0}^{n} t \cdot C_t\,(1+i)^{-t-1}}{\sum_{t=0}^{n} C_t\,(1+i)^{-t}} = \frac{\green{\bar{d}}}{1+i}
\]
\item When $i$ is a semi-annual rate: $\bar{v} = \dfrac{\bar{d}/2}{1+i}$
\medskip
\item When $i$ is an $m$-thly rate: $\bar{v} = \dfrac{\bar{d}/m}{1+i}$
\end{itemize}

\end{frame}

\begin{frame}{Approximating the change in price}

If we know the cash flows' \orange{volatility}, we can easily approximate the change in price due to a change in the effective yield, $i$:
\[
\purple{\%\Delta P(i)} \approx (\Delta i)\, \frac{\frac{d}{di} P(i)}{P(i)} = -\orange{\bar{v}}\,(\Delta i)
\]
\[
\purple{P(i + \Delta i)} \approx P(i) \cdot \left[1 - \orange{\bar{v}}\,(\Delta i)\right]
\]

\end{frame}

\begin{frame}{Example 9.4}

Consider the same bond as in Example 9.3: a 4-year bond with 6\% annual coupons priced at par.
\begin{enumerate}[a)]
\item Find the modified duration (volatility) of this bond.
\vspace{6em}
\item What is the approximate percentage change in the price of this bond in response to a 75 bps increase in the effective yield?
\end{enumerate}

\vfill
\end{frame}

\begin{frame}{Example 9.4, continued}

\begin{enumerate}[a)]
\setcounter{enumi}{2}
\item What is the \underline{exact} percentage change in the price of this bond in response to a 75 bps increase in the effective yield?

\vspace{8em}
How good was the approximation in part (b)?
\end{enumerate}

\vfill
\end{frame}

\begin{frame}{Portfolio duration}

\begin{itemize}
\item Suppose a portfolio consists of $n$ assets, with $n$ separate series of annual cash flows. Using an effective annual interest rate of $i$, asset $k$ has price $P_k(i)$ and duration $D_k$.
\medskip
\item The total price of the portfolio is $P(i) = \sum_{k=1}^{n} P_k(i)$ and the duration of the portfolio is
\end{itemize}
\[
D = -\frac{P'(i)}{P(i)} = -\frac{\sum_{k=1}^{n} P_k'(i)}{P(i)} = -\sum_{k=1}^{n} \frac{P_k'(i)}{P_k(i)} \frac{P_k(i)}{P(i)} = \sum_{k=1}^{n} D_k \frac{P_k(i)}{P(i)}
\]

\end{frame}

%======================================================================
\section{Convexity}
%======================================================================

\begin{frame}{9.3 Convexity}

\begin{itemize}
\item Modified duration works well to estimate the price change for small changes in $i$ ($\sim$less than 50 bps).
\medskip
\item For larger changes in yield, we need to take into account the \textcolor{red}{convexity} of the price function.
\end{itemize}

\begin{center}
\pricecurve{0.65}
\end{center}

\end{frame}

\begin{frame}{Second-order approximation}

\begin{itemize}
\item The price change resulting from a change in the annual effective yield is
\[
P(i + \Delta i) = P(i) + (\Delta i) P'(i) + \frac{(\Delta i)^2 P''(i)}{2!} + \cdots
\]
\item A second-order approximation gives
\[
P(i + \Delta i) \approx P(i) \left[1 - (\Delta i)\bar{v} + \frac{1}{2}(\Delta i)^2 \frac{P''(i)}{P(i)}\right]
\]
\end{itemize}

\end{frame}

\begin{frame}{Convexity}

Define the \textcolor{red}{convexity}, \textcolor{red}{$\bar{c}$}, of a set of cash flows as the \textcolor{red}{second derivative} of the price function with respect to the effective yield, divided by the price
\begin{align*}
P(i) = \sum_{t=0}^{n} C_t\,(1+i)^{-t} \;\Rightarrow\; \frac{d}{di} P(i) &= \sum_{t=0}^{n} -t \cdot C_t\,(1+i)^{-t-1} \\
\frac{d^2}{di^2} P(i) &= \sum_{t=0}^{n} t(t+1) \cdot C_t\,(1+i)^{-t-2} \\
\textcolor{red}{\bar{c}} = \frac{P''(i)}{P(i)} &= \frac{\sum_{t=0}^{n} t(t+1) \cdot C_t\,(1+i)^{-t}}{(1+i)^2 P(i)}
\end{align*}

\end{frame}

\begin{frame}{Example 9.5}

Consider the same bond as in Example 9.4: a 4-year bond with 6\% annual coupons, priced at par.
\begin{enumerate}[a)]
\item Find the convexity of this bond.
\end{enumerate}

\begin{center}
\renewcommand{\arraystretch}{1.3}
\begin{tabular}{ccc}
\toprule
$t$ & $C_t$ & $t(t+1) \cdot C_t (1+i)^{-(t+2)}$ \\
\midrule
1 & $0.06F$ & $1(2) \cdot 0.06F/(1.06)^3$ \\
2 & $0.06F$ & $2(3) \cdot 0.06F/(1.06)^4$ \\
3 & $0.06F$ & $3(4) \cdot 0.06F/(1.06)^5$ \\
4 & $1.06F$ & $4(5) \cdot 1.06F/(1.06)^6$ \\
\bottomrule
\end{tabular}
\end{center}

\vfill
\end{frame}

\begin{frame}{Example 9.5, continued}

\begin{enumerate}[a)]
\setcounter{enumi}{1}
\item What is the approximate percentage change in the price of this bond in response to a 75 bps increase in the effective yield, using convexity?

\vspace{8em}
How good is this approximation?
\end{enumerate}

\vfill
\end{frame}

\begin{frame}{Portfolio convexity}

The convexity of a portfolio is the \textcolor{red}{weighted average of the convexity of each item} in the portfolio, with the weight being the price of the item in proportion to the price of the entire portfolio:
\[
C = \sum_{k=1}^{n} \bar{c}_k \frac{P_k(i)}{P(i)}
\]

\end{frame}

%======================================================================
\section{Immunization}
%======================================================================

\begin{frame}{9.4 Immunization}

\begin{itemize}
\item Suppose we have a series of liability cashflows and we are trying to construct an asset portfolio so that the surplus (PV of assets $-$ PV of liabilities) does not become negative when interest rates change
\medskip
\item Two choices:
\begin{enumerate}
\item \purple{Redington immunization}
\item \green{Full immunization}
\end{enumerate}
\end{itemize}

\end{frame}

\begin{frame}{Redington immunization}

\begin{itemize}
\item Three conditions:
\begin{enumerate}
\item $PV_A(i) \ge PV_L(i)$
\item $\bar{d}_A(i) = \bar{d}_L(i)$
\item $\bar{c}_A(i) > \bar{c}_L(i)$
\end{enumerate}
\end{itemize}

\end{frame}

\begin{frame}{Redington immunization, cont'd}

\begin{itemize}
\item Ensures that $PV_A(i) \ge PV_L(i)$ for any \purple{small shift} in the yield rate
\end{itemize}

\begin{center}
\begin{tikzpicture}[scale=0.9]
  \draw[->, thick] (0,0) -- (7,0) node[right] {yield};
  \draw[->, thick] (0,0) -- (0,3.6) node[above left] {Price};
  % PV_L: black, PV_A: red dashed, tangent at x = 1.5
  \draw[thick, domain=0.6:4.6, samples=60] plot (\x, {0.6 + 1.5/\x});
  \draw[thick, red, dashed, domain=0.85:4.1, samples=60]
    plot (\x, {0.6 + 1.5/\x + 0.3*(\x-1.5)*(\x-1.5)/(1+0.8*(\x-1.5)*(\x-1.5))});
  \draw[densely dashed, gray] (1.5,0) -- (1.5,1.6);
  \node[below] at (1.5,-0.1) {$i$};
  \node[red, right] at (4.2,1.55) {$PV_A$};
  \node[right] at (4.6,0.85) {$PV_L$};
\end{tikzpicture}
\end{center}

\end{frame}

\begin{frame}{Assumptions (Redington imm'n)}

\begin{itemize}
\item Flat yield curve
\medskip
\item Parallel shifts in the yield curve (yields change by the same amount for all maturities)
\end{itemize}

\end{frame}

\begin{frame}{Example 9.6}

An insurance company has committed to making a payment of \$100,000 in 10 years.

\medskip
In order to fund its liability, the company has invested \$27,919.74 in a 5-year zero-coupon bond and \$27,919.74 in a 15-year zero-coupon bond.

\medskip
The annual effective yield is 6\%.

\medskip
Determine whether the company's position is immunized.

\end{frame}

\begin{frame}{Example 9.6, continued}

\begin{center}
\begin{tikzpicture}
  \node[anchor=west] at (-3,0.6) {\small Liability CF:};
  \draw (0,0) -- (7.5,0);
  \foreach \x/\lab in {0/0, 2.5/5, 5/10, 7.5/15} {
    \draw (\x,0.12) -- (\x,-0.12);
    \node[below] at (\x,-0.12) {\scriptsize $\lab$};
  }
  \node[anchor=west] at (-3,-1.0) {\small Asset CF:};
\end{tikzpicture}
\end{center}

\vfill
\end{frame}

\begin{frame}{Full immunization}

\begin{itemize}
\item A portfolio is \green{fully immunized at $i$} if
\[
S(i + \varepsilon) \ge 0 \text{ for certain } \varepsilon > 0
\]
\item Assuming a flat yield curve and parallel shifts in yields, a \orange{single cash flow} is \green{fully immunized} if
\begin{enumerate}
\item $PV_A(i) = PV_L(i)$
\item $\bar{d}_A(i) = \bar{d}_L(i)$
\item Asset cash flows occur both before and after the liability cash flow
\end{enumerate}
\end{itemize}

\end{frame}

\begin{frame}{Full immunization, continued}

\begin{center}
\begin{tikzpicture}
  \node[anchor=west] at (-3,0) {Cash flows:};
  \draw (0,0) -- (7,0);
  \foreach \x/\lab/\cf in {0/0/, 2.5/t-q/A_1, 4.5/t/L, 6.5/t+q/A_2} {
    \draw (\x,0.12) -- (\x,-0.12);
    \node[below] at (\x,-0.12) {$\lab$};
    \node[above] at (\x,0.15) {$\cf$};
  }
\end{tikzpicture}
\end{center}

where $PV_{A_1}(i) = PV_{A_2}(i)$.

\begin{center}
\renewcommand{\arraystretch}{1.3}
\begin{tabular}{lcc}
\toprule
 & \textbf{Liability} & \textbf{Asset portfolio} \\
\midrule
Duration & $t$ & $t$ \\
Convexity & $t(t+1)$ & $> t(t+1)$ \textcolor{red}{(why?)} \\
\bottomrule
\end{tabular}
\end{center}

\end{frame}

\begin{frame}{Full immunization, continued}

\begin{itemize}
\item To fully immunize a \orange{series of liability cash flows}, choose assets to \green{fully immunize each liability} cash flow separately.
\medskip
\item Some practical issues:
\begin{itemize}
\item Yield curves are usually not flat
\item Changes in yield curve are not parallel
\item Asset cash flows contain risk
\item Uncertainty in liability cash flows
\end{itemize}
\end{itemize}

\end{frame}

\begin{frame}{Example 9.7}

Liability payments of 100 each are due to be paid in 2, 4 and 6 years from now. Asset cash flow consists of $A_1$ in 1 year and $A_5$ in 5 years. The yield for all payments is 10\%.

\medskip
An attempt is made to have the asset cash flow immunize the liability cash flow by matching present value and duration.

\end{frame}

\begin{frame}{Example 9.7, continued}

\begin{enumerate}[a)]
\item Find $A_1$ and $A_5$.
\end{enumerate}

\vfill
\end{frame}

\begin{frame}{Example 9.7, continued}

\begin{enumerate}[a)]
\setcounter{enumi}{1}
\item Determine whether the conditions for Redington immunization are satisfied.
\end{enumerate}

\vfill
\end{frame}

\end{document}
